\documentclass{article}
\usepackage[utf8]{inputenc}
\usepackage{geometry}
\usepackage{fancyhdr}
\usepackage{listings}
\usepackage{xcolor}
\usepackage{framed}
\usepackage{tcolorbox}
\usepackage{amsmath}
\usepackage{amssymb}
\usepackage[T1]{fontenc}
\usepackage{caption}
\usepackage{babel}
\usepackage{hyperref}
\usepackage{titlesec}
\usepackage{graphicx}
\usepackage{color}
\usepackage[dvipsnames]{xcolor} 
\usepackage{media9}
\usepackage{hyperref}




\geometry{
  a4paper,
  left=25mm,
  right=25mm,
  top=25mm,
  bottom=25mm}

\hbadness=10000
\pagestyle{fancy}
\fancyhf{} 
\fancyfoot[C]{\thepage}
\fancyhead[L]{\textbf{MPI}} 
\fancyhead[R]{2024$-$2025$-$2026} 
\fancyhead[C]{\textbf{Alban COADIC}} 

\begin{document}
\begin{center}
\huge \textbf{TIPE:\@ Présentation TIPE 19--11--25}
\end{center}

\vspace{5mm}

\raggedright{\Large {\textcolor{OliveGreen}{\underline{\textbf{Sujet$ \ $:}}}}}\newline

\large Code en Rust d'une simulation de mécanique des fluides en 2D afin
de représenter les vortex de Von Karman.
\vspace{3mm}\\

\raggedright{\Large {\textcolor{OliveGreen}{\underline{\textbf{Objectif$ \ $:}}}}}
\vspace{3mm}\\

\large Comprendre les notions de base de la mécanique des fluides, les appliquer à un cas 
concret et étudier les potentielles méthodes de protection des batiments et structures face
aux vents (de manière plus générale face aux écoulements fluides).

\vspace{3mm}
\Large {\textcolor{OliveGreen}{\underline{ \textbf{Problématique$ \ $:}}}}
\vspace{3mm}\\

\large Comment modéliser les écoulements fluides autour d'un obstacle et étudier les
effets de ces écoulements sur les structures environnantes?\newline


\large 

\Large {\textcolor{OliveGreen}{\underline{\textbf{Bibliographie$ \ $:}}}}\newline

\large Cours de mécanique des fluides de l'université de Lyon$ \ $:

\begin{tcolorbox}[colback=white, colframe=gray]
  \normalsize\hypersetup{urlcolor=blue}
  \href{https://perso.univ-lyon1.fr/marc.buffat/COURS/BOOK_MECAFLU_HTML/intro.html}{Cours de mécanique des fluides - Université Lyon 1}
\end{tcolorbox}

\vspace{2em}

\large Thèse de BENSEDIRA Sidali de l'université de Blida$ \ $:

\begin{tcolorbox}[colback=white, colframe=gray]
  \normalsize\hypersetup{urlcolor=blue}
  \href{https://di.univ-blida.dz/jspui/bitstream/123456789/9666/1/32-530-642-1.pdf}{Thèse - Université de Blida}
\end{tcolorbox}

\vspace{2em}

\large Thèse de COUDERC Frédéric de l'université de Toulouse$ \ $:

\begin{tcolorbox}[colback=white, colframe=gray]
  \normalsize\hypersetup{urlcolor=blue}
  \href{https://theses.hal.science/tel-00143709}{Thèse - Université de Toulouse}
\end{tcolorbox}

\vspace{2em}

\begin{tcolorbox}[colback=lightgray!30, colframe=black]
\begin{center}
\hypersetup{pdfnewwindow=true}
\href{Maths et Notions.pdf}{\Large \textcolor{RoyalBlue}{\textbf{Maths et notions}}}
\end{center}
\end{tcolorbox}
\vspace{2mm}
\begin{tcolorbox}[colback=lightgray!30, colframe=black]
\begin{center}
\hypersetup{pdfnewwindow=true}
\href{Fancy_content_TIPE.pdf}{\Large \textcolor{RoyalBlue}{\textbf{Fancy content}}}
\end{center}
\end{tcolorbox}

\vspace{2mm}
\begin{tcolorbox}[colback=lightgray!30, colframe=black]
\begin{center}
\hypersetup{pdfnewwindow=true}
\href{Explication_code.pdf}{\Large \textcolor{RoyalBlue}{\textbf{Explication du code}}}
\end{center}
\end{tcolorbox}



\end{document} 